\PilotPage{5}
\section*{Pendulum}

\PilotFigure{p005-pendulum}{Pendulum geometry and free-body diagram.}

For the pendulum, let $u=L\theta$ denote the tangential displacement and let $W=mg$.
The sum of forces in the tangential direction is
\begin{equation}
    \sum F_u=-W\sin\theta.
    \tag{1}
\end{equation}
Newton's second law gives
\begin{equation}
    \sum F_u=m\ddot{u}.
    \tag{1a}
\end{equation}
Since $u=L\theta$,
\begin{equation}
    -W\sin\theta=mL\ddot{\theta}.
    \tag{2}
\end{equation}
For a small angle, $\sin\theta\approx\theta$ (radians), so
\begin{equation}
    -W\theta=mL\ddot{\theta},
    \tag{3}
\end{equation}
and therefore the equation of motion is
\begin{equation}
    mL\ddot{\theta}+W\theta=0.
    \tag{4}
\end{equation}
Using $W=mg$ in (4) gives
\begin{equation*}
    mL\ddot{\theta}+mg\theta=0,
\end{equation*}
or the governing equation
\begin{equation}
    \ddot{\theta}+\frac{g}{L}\theta=0.
    \tag{5}
\end{equation}

Assume a solution of the form
\begin{equation}
    \theta=e^{st}, \qquad \ddot{\theta}=s^2e^{st}=s^2\theta.
    \tag{6}
\end{equation}
Substitution into (5) gives
\begin{equation}
    \left(s^2+\frac{g}{L}\right)\theta=0,
    \tag{7}
\end{equation}
with characteristic equation
\begin{equation}
    s^2+\frac{g}{L}=0.
    \tag{8}
\end{equation}
The roots are
\begin{equation}
    s=\pm i\sqrt{\frac{g}{L}}.
    \tag{9}
\end{equation}
